\section{Finite bases: weights, duality, and the geometric comparison}

Throughout this section the variance is fixed by
$\mathbb D_{\mathcal C}(S)=\operatorname{Ho}\operatorname{Fun}(S^{\mathrm{op}},\mathcal C)$.
Thus $p_*^S$ is a limit over $S^{\mathrm{op}}$ and $p_!^S$ is a
colimit over $S^{\mathrm{op}}$. A homotopy finite category means a
category equivalent to a finite skeletal category with no nonidentity
endomorphisms. This is a condition on the indexing category, not merely
on the weak homotopy type of its nerve. In particular a nontrivial
finite group, regarded as a one-object category, is not homotopy finite.

\begin{claim}{G.1}{Known}{The exact stability input and the available tensor action}
In a stable derivator satisfying Der~1--4, homotopy finite limits
commute with arbitrary homotopy colimits. Every such derivator has a
canonical action by the derivator of spectra. Consequently a right
homotopy finite right Kan extension can be expressed as a weighted
homotopy colimit with a spectral profunctor as weight. No internal
monoidal structure on the coefficient derivator and no strongness
axiom are included in this assertion.
\end{claim}
\begin{proof}
The commutation theorem is \citet[Theorem 2.14]{DGGrothShulman2017}.
The spectral action is the consequence of
\citet[Theorem A.11 and Appendix A.3]{DGCisinskiTabuada2011}:
an object $X$ determines the cocontinuous morphism from spectra sending
the sphere to $X$. Their Definition A.1 separates ``strong'' from
``triangulated or stable''; Theorem A.11 is stated for any triangulated
derivator. The resulting coherent action, not an arbitrary choice of
tensor product, is the one used here. The weighted Kan-extension
statement is precisely \citet[Theorem 4.11(iv)]{DGGrothShulman2017},
applied to this module structure; their Construction 4.8 and Lemma 4.10
give the coend and its Yoneda identity. Their convention is covariant
diagrams, so it is applied to the opposite of the indexing functor
when using the presheaf convention of this paper.
\end{proof}

\begin{claim}{G.2}{Proved}{The canonical dualizing weight and its exact variance}
Let $S$ be homotopy finite. Write
\[
 R_s(t)=\Sigma^\infty_+\operatorname{Hom}_S(t,s),\qquad
 w_S(s)=\lim_{t\in S^{\mathrm{op}}}R_s(t).
\tag{G.2a}
\]
Then $w_S:S\to\mathrm{Sp}$ is a covariant spectral diagram, with
finite-spectrum values, and for every stable derivator there is a
canonical natural isomorphism
\[
 p_*^S F\simeq w_S\otimes_{[S]}F
       =\int^{s\in S} w_S(s)\otimes F(s).
\tag{G.2b}
\]
The integral denotes the derived coend, with its coherent spectral
action. It is not an objectwise tensor followed by the ordinary
colimit over $S^{\mathrm{op}}$.
\end{claim}
\begin{proof}
The spectral Yoneda identity gives the coherent density expression
\[
 F(t)\simeq\int^{s\in S}R_s(t)\otimes F(s).
\]
Apply the homotopy limit over $t\in S^{\mathrm{op}}$. By G.1 this
limit commutes with the homotopy colimits defining the coend. It also
commutes with the tensor action in its spectral variable: the action
preserves colimits, and a cocontinuous functor between stable
derivators preserves finite homotopy limits. The result is (G.2b).
This is also the terminal-target instance of the weight in
Groth--Shulman's Theorem 4.11(iv), so the calculation is coherent in
every external diagram variable, not just an objectwise calculation
in the underlying category.

Replace $S$ by a finite skeletal representative. Every hom-set is
finite, so every $R_s(t)$ is a finite wedge of sphere spectra. A limit
over this finite direct shape is a finite homotopy limit; finite
spectra are closed under those limits. This proves finiteness of
$w_S(s)$. Equivalent indexing categories give canonically equivalent
weights after transport. For $S=\varnothing$, the coend and the limit
are both the zero object in a stable category.
\end{proof}

\begin{claim}{G.3}{Refuted}{A pointwise dualizing-object twist already fails on one arrow}
The assertion that every homotopy finite $S$ admits
$\omega\in\operatorname{Fun}(S^{\mathrm{op}},\mathrm{Sp})$ with
\[
 p_*^SF\simeq p_!^S(\omega\otimes F)
 \quad\hbox{for all }F
\tag{G.3}
\]
where the tensor is pointwise, is false. It remains false for
pointwise perfect complexes over a nonzero field.
\end{claim}
\begin{proof}
Take $S=[1]=(0\to1)$. A presheaf is an arrow $F_1\to F_0$.
The initial object of $S^{\mathrm{op}}$ is $1$, and its terminal
object is $0$; therefore $p_*F=F_1$ and $p_!F=F_0$.
Choose $F=(M\to0)$, with $M$ any nonzero finite spectrum, or the
field concentrated in degree zero. The right side of (G.3) is
$\omega_0\otimes0=0$, while the left side is $M$. This contradiction
does not depend on the value, dualizability, or invertibility of
$\omega$. In contrast, the actual weight in (G.2a) is
$w_{[1]}=(0\to\mathbb S)$, a diagram on $S$, and its coend with
$F$ is $F_1$. The endpoint and weight calculations are included in
\texttt{code/dg\_finite\_checks.py}.
\end{proof}

\begin{claim}{G.4}{Proved}{Transitivity and base change belong to the kernels}
For a functor $u:R\to S$ between homotopy finite categories, set
\[
 K_u(s,r)=(u_*R_r)(s),\qquad
 K_u\in\operatorname{Fun}(S^{\mathrm{op}}\times R,\mathrm{Sp}).
\]
Then $u_*F\simeq K_u\otimes_{[R]}F$. For $v:S\to T$ there are
canonical composition and transitivity identities
\[
 K_{vu}(t,r)\simeq\int^{s\in S}K_v(t,s)\wedge K_u(s,r),\qquad
 w_R(r)\simeq\int^{s\in S}w_S(s)\wedge K_u(s,r).
\tag{G.4}
\]
A homotopy exact square gives the corresponding base-change
equivalence of these kernels and of their weighted functors. No
claim that every cartesian square is homotopy exact is needed.
\end{claim}
\begin{proof}
The comma categories computing $u_*$ are finite and have no
nonidentity endomorphisms after replacing the indexing categories
by finite skeletal representatives. Hence $u_*$ preserves homotopy
colimits by G.1. Apply $u_*$ to the density expression for $F$ in
G.2. This gives the first formula and proves its naturality.
Apply the formula twice to $v_*u_*F$, and use Fubini for homotopy
coends. The canonical isomorphism $(vu)_*\simeq v_*u_*$ identifies
the resulting kernel, proving the first identity in (G.4).
Set $T=e$ to obtain the second identity. Associativity of these
isomorphisms is the associativity of Kan extension and of coends;
therefore this is a coherent transitivity formula.

For a square whose specified Beck--Chevalley mate is invertible,
evaluate that mate on all the spectral representables $R_r$.
This is the asserted kernel comparison. Conversely, when the
functors concerned preserve colimits, their density expressions
show that equivalence on these representables gives the equivalence
on every diagram. Thus the compatibility is with the actual
base-change mate and its pasted composites. This argument does not
replace that mate by an unrelated isomorphism of functors.
\end{proof}

\begin{claim}{G.5}{Refuted}{The weight is not invariant under unrestricted restriction}
The formula $i^*w_S\simeq w_R$ for every functor $i:R\to S$ is
false, even for finite posets. Consequently no base-change statement
for a ``dualizing object'' may be supplied by ordinary restriction
without an additional hypothesis or an extraordinary pullback.
\end{claim}
\begin{proof}
Let $i:e\to[1]$ select $0$. G.3 computes
$i^*w_{[1]}=0$, whereas $w_e=\mathbb S$. The nonzero sphere cannot
be equivalent to zero. The same example shows that the transform
in G.9 below does not commute with this restriction: its value at
$0$ is always zero, whereas the transform on $e$ is the identity.
The kernel transitivity formula G.4 is the valid replacement.
\end{proof}

\begin{claim}{G.6}{Proved}{A coefficient duality on precisely specified finite objects}
For every small $S$, pointwise Spanier--Whitehead duality gives
\[
 \operatorname{Fun}(S^{\mathrm{op}},\mathrm{Sp}^{\mathrm{fin}})^{\mathrm{op}}
 \simeq\operatorname{Fun}(S,\mathrm{Sp}^{\mathrm{fin}}).
\tag{G.6a}
\]
For a ring $A$, derived $A$-linear duality similarly gives
\[
 \operatorname{Fun}(S^{\mathrm{op}},\operatorname{Perf}(A))^{\mathrm{op}}
 \simeq\operatorname{Fun}(S,\operatorname{Perf}(A^{\mathrm{op}})).
\tag{G.6b}
\]
For homotopy finite $S$, all limits and colimits in these formulas
remain perfect and
$D(p_*^SF)\simeq p_!^{S^{\mathrm{op}}}(DF)$.
This is a contravariant duality. In (G.6b) the coefficient ring can
be kept unchanged when it is commutative.
\end{claim}
\begin{proof}
For a finite spectrum $E$, the evaluation map
$E\to\operatorname{Map}(\operatorname{Map}(E,\mathbb S),\mathbb S)$
is an equivalence: it holds for the sphere and is preserved by finite
cofibres, suspensions, and retracts. This proves the equivalence on
finite spectra. Applying it to values and reversing structure maps
gives (G.6a), including all coherences. For perfect $A$-modules the
same proof starts from finite free modules and uses
$\operatorname{RHom}_A(-,A)$; the dual is a module on the opposite
side, which accounts for $A^{\mathrm{op}}$.

A limit over a homotopy finite shape is constructed by finite
homotopy limits, and finite spectra and perfect modules are closed
under those operations. An equivalence to an opposite category
exchanges their limits and colimits. Applying this to the diagram
gives the last formula. No assertion is made that arbitrary
spectra or arbitrary modules are reflexive under this duality.
\end{proof}

\begin{claim}{G.7}{Refuted}{Finite EI is not a coefficient-independent extension of homotopy finite}
The weighted-colimit assertion for all stable coefficients does not
extend from homotopy finite categories to all finite EI categories.
For $S=BC_2$ and coefficients $k=\mathbb F_2$, the functor $p_*^S$
does not preserve countable coproducts, and hence cannot be any
fixed weighted homotopy colimit on all diagrams.
\end{claim}
\begin{proof}
The category has one object and exactly the arrows $1,g$, with
$g^2=1$. It is finite and every endomorphism is invertible. Write
$A=k[C_2]=k[\epsilon]/(\epsilon^2)$, where $\epsilon=g-1$.
The trivial module $k$ has the bounded-above projective resolution
\[
 \cdots\xrightarrow{\epsilon}A\xrightarrow{\epsilon}A
 \xrightarrow{\epsilon}A\longrightarrow k\longrightarrow0,
\]
with $A$ in every degree $0,-1,-2,\ldots$. Exactness follows from
$\ker(\epsilon:A\to A)=\epsilon A=\operatorname{im}(\epsilon)$
and $\ker(A\to k)=\epsilon A$. A bounded-above complex of
projectives computes derived Hom into any complex.

Let $M_n=k[n]$ for $n\geq0$ in cohomological notation, with
trivial action, and let $M=\bigoplus_{n\geq0}M_n$.
In degree zero the Hom complex from the displayed resolution to
$M$ is $\prod_{n\geq0}k$. Every differential is zero, since
$\epsilon$ acts trivially on the target. Thus
\[
 H^0\operatorname{RHom}_A(k,M)=\prod_{n\geq0}k,
 \qquad
 \bigoplus_{n\geq0}H^0\operatorname{RHom}_A(k,M_n)
       =\bigoplus_{n\geq0}k.
\]
The comparison is the usual inclusion of the direct sum into the
product and omits the sequence $(1,1,\ldots)$. Since
$p_*^S=\operatorname{RHom}_A(k,-)$, it fails to preserve this
coproduct. A fixed weighted coend preserves all colimits in its
input, so none can represent this functor. The multiplication table,
associativity, and the resolution's differential matrices are checked
in \texttt{code/dg\_finite\_checks.py}; the infinite product argument
above, rather than a finite computation, proves the obstruction.
\end{proof}

\begin{claim}{G.8}{Proved}{A finite-group case with its necessary coefficient hypothesis}
Let $G$ be a finite group and $k$ a commutative ring in which
$|G|$ is invertible. For the derivator of $k$-modules and
$S=BG$, the norm gives a natural equivalence $p_!^S\simeq p_*^S$
on all unbounded complexes. In particular this finite EI case has
ambidexterity, although G.7 rules out a uniform assertion without
the coefficient hypothesis.
\end{claim}
\begin{proof}
For any $k[G]$-module $M$, the map
$[m]\mapsto\sum_{g\in G}gm$ from coinvariants to invariants
has inverse $x\mapsto |G|^{-1}[x]$. The averaging idempotent
$|G|^{-1}\sum_g g$ makes invariants an exact direct summand
functor. The trivial $k[G]$-module is projective on either side,
so coinvariants and invariants compute their derived functors
degree by degree on all complexes. The displayed natural inverse
therefore supplies the derived norm equivalence. This is a proved
class of finite EI examples, not a classification of all finite EI
categories or all coefficient systems.
\end{proof}

\begin{claim}{G.9}{Known}{The exact canonical equivalence criterion for finite posets}
For a finite poset $P$ define the colimit-preserving transform
\[
 (V_PF)(s)=
 \lim_{t\in P^{\mathrm{op}}}
  \bigl(\Sigma^\infty_+\operatorname{Hom}_P(t,s)\otimes F(t)\bigr),
 \qquad V_P:\mathbb D(P)\longrightarrow\mathbb D(P^{\mathrm{op}}).
\tag{G.9a}
\]
There is always a canonical formula
\[
 p_*^P\simeq p_!^{P^{\mathrm{op}}}V_P.
\tag{G.9b}
\]
For spectra the transform (G.9a) is an equivalence if and only if
each strict upper interval $P_{>p}$ is Gorenstein* over $\mathbb Z$.
Here a finite poset $L$ is Gorenstein* when every open interval in
$L$ with a newly adjoined bottom and top has the integral homology
of a sphere; the empty interval is allowed as the $(-1)$-sphere.
When this condition holds, (G.9a) gives an equivalence for every
stable derivator, and it restricts to pointwise finite spectra or
perfect complexes. This is a condition on this canonical transform;
it is not a classification of every possible equivalence between
the two diagram categories.
\end{claim}
\begin{proof}
Apply \citet[Theorem A and Proposition 5.12]{DGAoki2022} to
$Q=P^{\mathrm{op}}$. His diagram category is
$\operatorname{Fun}(Q,\mathrm{Sp})$, so his lower-interval
condition becomes the upper-interval condition just stated.
His Definition 3.9 and Proposition 3.11 give the precise
Gorenstein* convention, including the empty case. This proves the
if-and-only-if assertion with its exact variance.

For completeness, formula (G.9b) itself follows without that
condition. Commute the finite limit over $t$ with the colimit over
$s$ in (G.9a). For fixed $t$, the covariant diagram
$s\mapsto\operatorname{Hom}_P(t,s)$ has category of elements
$t/P$, which has an initial object. Its homotopy colimit is
contractible, so the colimit of its suspension spectra is the
sphere. The result is $\lim_t F(t)$, as required.

The extension of the equivalence to stable derivators can be made
without an enhancement assumption on that derivator. By G.1 and
the density argument of G.4, $V_P$ is tensoring with a spectral
bimodule. An inverse to its colimit-preserving equivalence for
spectral diagrams is also colimit preserving, and density represents
that inverse by a spectral bimodule. The two composites are the
identity bimodules: evaluate the two equivalences on representables.
Tensoring these bimodule identities with any stable derivator
proves the assertion. For a finite direct shape, an objectwise
perfect diagram is compact, by its finite normalized bar resolution
in representables, and every compact diagram is objectwise perfect.
The equivalence preserves compact objects; hence so does its
inverse. This proves the stated restriction to finite objects.
\end{proof}

\begin{claim}{G.10}{Proved}{Verification of the finite least-vertex geometric dictionary}
Let $P$ be a finite poset and $Z=|NP|$. Label a point by the
least vertex with positive barycentric coordinate. With downsets
open on $P$, this is a continuous map $s:Z\to P$; its strata are
contractible, the stratification is conical, and its exit-path
infinity-category is $P^{\mathrm{op}}$. Consequently
\[
 \operatorname{Shv}_{P\text{-}\mathrm{cons}}(Z;\mathcal C)
       \simeq\operatorname{Fun}(P^{\mathrm{op}},\mathcal C)
\tag{G.10}
\]
for spaces, for every compactly assembled coefficient
infinity-category, and for every stable presentable coefficient
infinity-category. The inverse-image functor from sheaves on the
finite downset Alexandrov space $P$ is the fully faithful functor
onto the subcategory in (G.10).
\end{claim}
\begin{proof}
For a downset $U$, the complement of $s^{-1}U$ is the closed
subcomplex on $P\setminus U$: a chain has minimum outside $U$
exactly when all its vertices lie outside $U$. This proves
continuity. The stratum labelled $p$ consists of chains with
minimum $p$ and positive coordinate at $p$; it is the open cone
on $|N(P_{>p})|$ and is contractible.

Here is an explicit local chart, which also checks the geometric
hypothesis instead of assuming it from the attached STRATA package.
The open star $t_p>0$ of a vertex $p$ has the homeomorphism
\[
 \{t_p>0\}\cong
 C_\infty|N(P_{<p})|\ \times\ C_\infty|N(P_{>p})|.
\tag{G.10a}
\]
The two radial coordinates are
$\sum_{q<p}t_q/t_p$ and $\sum_{q>p}t_q/t_p$; the remaining
coordinates are normalized on their respective chains. The inverse
divides all coordinates by $1+a+b$. If a radial coordinate is zero,
its normalized coordinates are collapsed to the cone vertex, so
these formulas extend continuously there. The second factor is a
factor inside the $p$-stratum. The first factor is the normal cone,
whose nonvertex strata have their original minimum labels below
$p$. Taking neighbourhoods in the second factor gives conical charts
at every point of the $p$-stratum. These are finite polyhedral
charts, hence locally singular and paracompact.

Refine this stratification to the open-simplex stratification of the
finite order complex. Its exit category is the face poset
$\operatorname{Ch}(P)$ of nonempty chains, ordered by inclusion,
by \citet[Theorem A.6.10]{DGLurieHA}. The coarsening map is
$m:\operatorname{Ch}(P)\to P^{\mathrm{op}}$,
$m(\sigma)=\min\sigma$. The coarsening theorem
\citet[Theorem 3.3.5]{DGHPT2026} identifies the new exit category
with the localization of $\operatorname{Ch}(P)$ at inclusions
preserving the minimum. Its hypotheses hold for this finite
conically stratified polyhedron, by
\citet[Theorems A.6.4 and A.9.3]{DGLurieHA}.
Finally $N\operatorname{Ch}(P)=\operatorname{Sd}(NP^{\mathrm{op}})$,
and $Nm$ is the last-vertex map for $P^{\mathrm{op}}$.
\citet[Proposition 6.3.6.2]{DGKerodonSubdivision} says exactly that
this map is universally localizing at those edges. Thus the
localized exit category is $P^{\mathrm{op}}$.

Space-valued exodromy is now HA, Theorem A.9.3. The coefficient
extension in the precise two classes stated above is
\citet[Corollary 4.1.15]{DGHPT2026}; the finite indexing poset is
noetherian, as required for its stable-coefficient case. Finally
sheaves on a finite downset Alexandrov space are diagrams on
$P^{\mathrm{op}}$, by finite gluing on principal opens, and the
map on exits induced by $s$ is the just computed equivalence.
Naturality of exodromy therefore identifies $s^*$ with (G.10).
This verifies the specific dictionary needed here and does not
assert its arbitrary-presentable-coefficient extension.
\end{proof}

\begin{claim}{G.11}{Refuted}{Minimum constructibility is not automatically preserved by Verdier duality}
Even for $P=[1]$, usual contravariant Verdier duality on $|NP|$
does not preserve the least-vertex constructible sheaves with
perfect stalks. Also, the canonical covariant transform $V_P$ in
G.9 is not an equivalence for this $P$.
\end{claim}
\begin{proof}
The strata of $[0,1]$ are $[0,1)$ and $\{1\}$. The constant
sheaf $k$ is constructible for this stratification. Its Verdier
dual is the dualizing sheaf of the closed interval. At the endpoint
$0$ the local costalk of $k$ is
\[
 \operatorname{fib}\bigl(R\Gamma([0,\epsilon);k)
             \to R\Gamma((0,\epsilon);k)\bigr)=0.
\]
At an interior point the punctured interval has two components,
and the corresponding restriction is the diagonal $k\to k\oplus k$.
Its fibre is $k[-1]$. Dualizing these finite costalks shows that
the stalk of the dualizing sheaf is zero at the endpoint and
$k[1]$ at an interior point; the other endpoint is the same as
$0$. Equivalently it is $j_!k_{(0,1)}[1]$.
It is not locally constant on $[0,1)$, so the proposed preservation
fails even for a finite, compact, contractible realization.
The local relative-cochain ranks are checked in
\texttt{code/dg\_finite\_checks.py}.

Algebraically, for $F=(F_1\to F_0)$, (G.9a) gives
$V_{[1]}F=(0\to F_1)$. The nonzero diagram $(0\to k)$ is sent
to zero. This agrees with Aoki's obstruction: the strict upper
interval of $0$ is a singleton, whose reduced homology is zero
and is not the homology of any sphere. This refutes the claimed
canonical identification on every homotopy finite base. It does
not refute an equivalence obtained by other operations: $[1]$ is
isomorphic to its opposite, and G.14 gives further such examples.
\end{proof}

\begin{claim}{G.12}{Known}{A correct geometric Verdier identification with a checked hypothesis}
Suppose the finite poset $P$ satisfies the upper-interval criterion
of G.9. For $Z=|NP|$ and its minimum map $s:Z\to P$, the
spectral transform has the identification
\[
 V_P\simeq s_+\,\mathbb V_Z\,s^*.
\tag{G.12}
\]
Here $\mathbb V_Z$ is Lurie's covariant sheaf--cosheaf duality and
$s_+$ is cosheaf pushforward, with value on $U$ given by the value
on $s^{-1}U$. The same identification holds for derived coefficients
over a ring. Thus it includes an explicit pushforward of cosheaves;
it does not require the ordinary sheaf-valued Verdier dual to stay
constructible for the same minimum strata.

A concrete class satisfying the criterion is
$P=\operatorname{Face}(K)^{\mathrm{op}}$ for any finite simplicial
complex $K$, with nonempty faces. In this case the minimum strata
on the barycentric subdivision are precisely the open simplices of
$K$, so the identification is the usual cellular geometric one.
\end{claim}
\begin{proof}
Use \citet[Theorem D, or Theorem 6.13 with Remark 6.14]{DGAoki2022}
with $Q=P^{\mathrm{op}}$. Its hypotheses are: $Q$ is Verdier,
the source topos is proper and separated, and the spectral
inverse-image functor is fully faithful. G.9 verifies the first;
the compact Hausdorff finite polyhedron $Z$ verifies the second;
G.10 verifies the third. Aoki uses the upsets-open convention
on $Q$, which is the downsets-open convention on $P$. His
identification is exactly (G.12), with these variances. The
ring-coefficient assertion follows by tensoring the colimit-preserving
spectral constructions with the module category; equivalently the
same spectral kernels act on that stable coefficient category.

For the stated simplicial example and a simplex $\sigma$, the
strict upper interval at $\sigma$ in $P$ is the opposite of the
proper nonempty face poset of $\sigma$. Its order complex is the
barycentric subdivision of the boundary sphere of $\sigma$.
All augmented open intervals are boundaries of simplex faces,
so the precise Gorenstein* condition holds, including an empty
interval at a vertex. The usual barycentric homeomorphism
$|N\operatorname{Face}(K)|\cong|K|$ sends a point with maximal
face $\sigma$ in its supporting chain to the relative interior
of $\sigma$: its coefficient on the barycenter of $\sigma$ is
positive, giving positive coordinate at every vertex of $\sigma$.
Conversely every point of that relative interior has that maximal
face in its supporting barycentric simplex. Minimum in the
opposite face poset is this maximum, proving the asserted match
of strata.
\end{proof}

\begin{claim}{G.13}{Refuted}{The weighted-limit argument does not persist unstably}
There is no fixed-colimit-weight description of all finite limits
of spaces. Consequently neither the stability argument of G.1
nor its duality conclusions may be exported to unstable fine
theories merely because the base is finite.
\end{claim}
\begin{proof}
Let $S$ have two objects and only identity arrows. Its limit is
the product. Put $F=(*,\varnothing)$ and
$G=(\varnothing,*)$. Then
\[
 (\lim F)\amalg(\lim G)=\varnothing,
 \qquad \lim(F\amalg G)=*.
\]
Thus this finite limit functor fails to preserve a binary
coproduct. A weighted colimit with fixed weight preserves
coproducts of the input, so it cannot be this functor. All four
values are finite discrete spaces and their cardinalities are
recorded by the finite-check program. Independently, D.3 gives
a finite poset whose two unstable fine categories are not even
equivalent. For a discrete finite base the two diagram categories
do coincide by opposition, but their product and coproduct
operations still differ as this example shows.
\end{proof}

\begin{claim}{G.14}{Known}{Tilting gives additional equivalences with different data}
If $T$ is the path category of a finite oriented tree, then
$\mathbb D(T)\simeq\mathbb D(T^{\mathrm{op}})$ for every stable
derivator, pseudonaturally in exact morphisms of derivators.
More generally any two orientations of the same finite tree
are strongly stably equivalent. These are reflection-functor
equivalences and are not the assertion that the transform
$V_T$ in G.9 is always an equivalence.
\end{claim}
\begin{proof}
This is \citet[Theorem 9.19 and Corollary 9.23]{DGGrothStovicek2016}.
The supplied source \texttt{tilting2.tex} was checked at the
reflection theorem and its tree corollary. A source reflection
replaces the value there by the cofiber of the map to the finite
biproduct of its outgoing neighbours; the inverse uses a fibre.
The published proof constructs these as coherent finite Kan
extensions and proves inverse equivalences. Successive source
and sink reflections realize any reorientation of a tree,
including complete reversal. Only finite limits and colimits
occur, which explains the asserted compatibility with exact
morphisms of stable derivators and, for spectra or rings,
preservation of pointwise perfect diagrams. These operations
depend on the chosen reflection sequence. In particular their
existence does not supply a dimension, an orientation, or a
pointwise dualizing-object formula such as the refuted (G.3).
\end{proof}
